\documentclass[11pt,letterpaper]{article}
\usepackage[T1]{fontenc}
\usepackage{lmodern}
\usepackage[left=.97in,right=.97in,top=.89in,bottom=.94in]{geometry}
\usepackage{amsmath,amssymb,amsthm,mathtools,microtype}
\usepackage{needspace}
\usepackage{xcolor}
\usepackage[colorlinks=true,linkcolor=blue!40!black,citecolor=blue!40!black,urlcolor=blue!40!black]{hyperref}
\hypersetup{pdftitle={A weighted reciprocal Lucas sum lies outside Q(sqrt(5))},
  pdfauthor={twinheterodox},pdfsubject={Irrationality of a weighted reciprocal Lucas sum}}
\urlstyle{same}
\newcommand{\Q}{\mathbb Q}
\newcommand{\Z}{\mathbb Z}
\newcommand{\C}{\mathbb C}
\newcommand{\LL}{\mathcal L}
\newcommand{\FF}{\mathcal F}
\newcommand{\ord}{\operatorname{ord}}
\newcommand{\logp}{\log^{+}}
\newtheorem{theorem}{Theorem}
\newtheorem{corollary}[theorem]{Corollary}
\newtheorem{lemma}[theorem]{Lemma}
\title{\Large A weighted reciprocal Lucas sum lies outside \(\Q(\sqrt5)\)}
\author{twinheterodox}
\date{18 November 2026}
\begin{document}
\maketitle
\begin{abstract}
For the even-indexed Lucas numbers \(L_{2n}\), the reciprocal sums with
numerators \(1\) and \(n^2\) are known to be transcendental. We prove
that the intermediate sum \(\sum_{n\ge1}n/L_{2n}\) lies outside
\(\Q(\sqrt5)\), and that the same conclusion holds with \(L_{2an}\) in
the denominator for every integer \(a\ge1\). More generally, the
weight-two value \(\beta_{1/b}(2)=\sum_{n\ge1}n/(b^n+b^{-n})\) of the
\(q\)-analogue of Dirichlet's beta function lies outside \(\Q(b)\) for
every real algebraic integer \(b>1\) whose other conjugates have modulus
at most one.
\end{abstract}

\section{Three weights for the even Lucas numbers}
The Lucas numbers satisfy \(L_0=2\), \(L_1=1\) and
\(L_{n+2}=L_{n+1}+L_n\). With the same denominators, consider the three sums
\[
 \sum_{n\ge1}\frac{1}{L_{2n}},\qquad
 \sum_{n\ge1}\frac{n}{L_{2n}},\qquad
 \sum_{n\ge1}\frac{n^2}{L_{2n}}.
\]
Duverney, Ke.~Nishioka, Ku.~Nishioka and Shiokawa proved that the first
sum is transcendental \cite[p.~97 of the Kyoto account]{dnns}. The
third is also transcendental, by the odd-weight result of Jouhet and
Mosaki \cite[\S1, (1.1) and the discussion before Theorem~1.1]{jm}.
What happens for the intermediate weight \(n\)? We prove that
\(1/3+2/7+3/18+4/47+\cdots\) lies outside \(\Q(\sqrt5)\), the field
generated by the roots of the Lucas recurrence.

To see the connection with beta values, write \(\varphi=(1+\sqrt5)/2\).
Binet's formula \(L_j=\varphi^j+(-\varphi^{-1})^j\) gives
\(L_{2n}=\varphi^{2n}+\varphi^{-2n}\).
For \(0<q<1\) and positive integers \(s\), Jouhet and Mosaki define
\[
 \beta_q(s)=\sum_{n\ge1}\frac{n^{s-1}q^n}{1+q^{2n}} .
\]
At \(q=\varphi^{-2}\) the three sums are precisely
\(\beta_q(1),\beta_q(2),\beta_q(3)\). The following theorem treats
the middle weight.

\begin{theorem}\label{thm:bases}
Let \(b>1\) be a real algebraic integer. If every conjugate of \(b\)
other than \(b\) has modulus at most one, then
\[
       \beta_{1/b}(2)=\sum_{n\ge1}\frac{n}{b^n+b^{-n}}\notin\Q(b).
\]
\end{theorem}

\begin{corollary}\label{cor:lucas}
For every integer \(a\ge1\), the sum
\(\displaystyle\sum_{n\ge1}n/L_{2an}\) lies outside \(\Q(\sqrt5)\).
\end{corollary}
\begin{proof}
Take \(b=\varphi^{2a}\). This is an algebraic integer with conjugate
\(\varphi^{-2a}<1\). These two values are distinct, so \(b\) generates
\(\Q(\sqrt5)\). Now \(b^n+b^{-n}=L_{2an}\), and Theorem~\ref{thm:bases}
applies.
\end{proof}

The base condition includes Pisot and Salem numbers. Bundschuh and
Zudilin singled out the arithmetic of this \(q\)-analogue of Catalan's
constant at algebraic \(q\), writing it as the alternating series in
Section~\ref{sec:pencil} \cite[\S6, p.~121]{bz}. For an integer
\(b\ge2\), Jouhet and Mosaki proved that at least one of the ten values
\(\beta_{1/b}(2),\beta_{1/b}(4),\ldots,\beta_{1/b}(20)\) is irrational
\cite[Corollary~1.3]{jm}. The integer-base theorem of uplift-universal
isolates \(\beta_{1/b}(2)\) \cite[Theorem~1.1]{uplift}. For the Lucas
denominators, Wang also relates the sums with even numerator powers to
derivatives of Eisenstein series \cite[Lemma~5.1 and Corollary~5.2 of the preprint]{wang}.

The Lucas example illustrates the obstacle at the other embeddings.
Every partial sum of \(\sum n/L_{2n}\) is rational. Yet if its limit
belonged to \(\Q(\sqrt5)\), the partial sums would give no formula for
its conjugate \(Y\). Indeed, conjugation sends \(\varphi^{-2m}\), which tends to zero, to
\(\varphi^{2m}\), which grows. We construct integer polynomials that
become small at \((b,\beta_{1/b}(2))\). At every other conjugate \(b'\)
we bound the same polynomials at \((b',Y)\) for every fixed \(Y\), even
when \(\lvert b'\rvert=1\). The product over the embeddings then gives
the contradiction to the lower bound for a nonzero integer norm.

We use the signed determinant and the cyclotomic divisibility of the
integer-base paper \cite[\S\S2--7]{uplift}. To compare conjugate
specializations, we use the base as a polynomial variable and bound the
normalized determinant throughout the closed unit disk. This follows the
conjugate estimate for a weight-three Lambert sum
\cite[\S\S4 and 6]{weight3}. Quadratic norms and conjugate determinant
bounds also appear in Krisman's proof for reciprocal Fibonacci cubes
\cite[\S5]{krisman}. For an earlier use of estimates over all places
for Lambert series at algebraic bases, see Bundschuh and Väänänen
\cite[\S2]{bv}. The proof below uses partial fractions,
elementary facts about algebraic integers and the maximum modulus principle.

\section{One determinant for all embeddings}\label{sec:pencil}
Fix a real \(b>1\) and set \(q=1/b\) and \(\xi=\beta_q(2)\).
The sum converges since its \(n\)-th term is at most \(nq^n\).
To obtain an integer polynomial in the unknown value, we need several
expressions involving only that value and rational functions of the base.
Expand the denominator in the positive series for \(\beta_q(2)\). With
\(\chi(2j)=0\) and \(\chi(2j+1)=(-1)^j\), this gives
\[
 \frac{q^j}{1+q^{2j}}=q^j-q^{3j}+q^{5j}-\cdots
                      =\sum_{m\ge1}\chi(m)q^{mj}\quad(j\ge1).
\]
Write \(g(y)=y/(1-y)^2\) for \(y\ne1\). When \(|y|<1\), its power
series is \(g(y)=\sum_{j\ge1}j y^j\).
Summing gives the identity \(\xi=\sum_{k\ge0}(-1)^k g(q^{2k+1})\) of
Jouhet and Mosaki \cite[\S6.2]{jm}. The interchange is absolute since
\(\sum_{m\ge1}q^m/(1-q^m)^2<\infty\).
At the Lucas parameter \(q=\varphi^{-2}\), Binet's formula gives
\(g(q^m)=(\varphi^m-\varphi^{-m})^{-2}=1/L_m^2\) for odd \(m\).
This identifies the two Lucas expansions
\begin{equation}\label{eq:odd-squares}
 \sum_{n\ge1}\frac{n}{L_{2n}}
 =\sum_{k\ge0}\frac{(-1)^k}{L_{2k+1}^{\,2}}
 =1-\frac{1}{4^2}+\frac{1}{11^2}-\frac{1}{29^2}+\cdots .
\end{equation}

For general \(q\), the odd exponents give the nodes \(q^m\), and applying
\(x\,d/dx\) to \(x^j\) supplies the numerator \(j\). Define, for
\(f\in C^1[0,q]\),
\[
  L(f)=\sum_{m\ge1}\chi(m)q^m f'(q^m).
\]
The sum is bounded in modulus by \(q\|f'\|_\infty/(1-q^2)\).

\Needspace{5\baselineskip}
Removing the first odd node changes the sign of the remaining series.
For example,
\(L((1-q^2x)^{-1})=g(q)-\xi\).
Removing the first \(a\) odd nodes gives the same value with a finite
correction and sign \((-1)^a\). For a variable base \(t\), record the
finite corrections and polynomial moments by
\begin{equation}\label{eq:moments}
 S_N(t)=\sum_{m=1}^N\chi(m)\frac{t^m}{(t^m-1)^2},\qquad
 \mu_0(t)=0,\qquad \mu_j(t)=\frac{jt^j}{1+t^{2j}}\quad(j\ge1).
\end{equation}
Indeed \(g(q^m)=g(b^m)\). Summing the geometric progression on odd
indices gives \(L(x^j)=\mu_j(b)\), and the shift just described gives
\[
 L\big((1-x/b^{2a})^{-1}\big)=(-1)^a\big(\xi-S_{2a}(b)\big).
\]

For size \(K\ge1\) write \(r_a=t^{2a}\) and
\(D_K(x)=\prod_{a=1}^K(1-x/r_a)\). On rational functions of \(x\) with
denominator dividing \(D_K\), define the linear functional \(\LL_{t,X}\)
over \(\Q(t,X)\) by
\[
 \LL_{t,X}(x^j)=\mu_j(t),\qquad
 \LL_{t,X}\big((1-x/r_a)^{-1}\big)=(-1)^a\big(X-S_{2a}(t)\big).
\]
These values specify the functional by polynomial division and simple
partial fractions. At \((t,X)=(b,\xi)\) it agrees with \(L\).
The determinant from \cite[\S2]{uplift} is
\begin{equation}\label{eq:delta}
 \Delta_K(t,X)=
 \det\!\left[\LL_{t,X}\!\left(\frac{x^{i+j}}{D_K(x)}\right)\right]_{0\le i,j<K}.
\end{equation}
Each entry is affine in \(X\), so the degree in \(X\) is at most \(K\).
At \(K=1\) the determinant is simply \(t/(t-1)^2-X\). At \(K=3\) the
bottom right entry involves \(x^4/D_3(x)\), whose polynomial quotient
is linear. The moments in \eqref{eq:moments} include its contribution.
At this size the two computations, by polynomial division and by the
pole recurrence, can be compared exactly at both Lucas embeddings. From the
root of the source bundle, run \texttt{python3}
\nolinkurl{checks/check_three_poles.py}.
Rational modifications of moments occur in Uvarov's transformation of
orthogonal polynomials \cite[\S1]{uvarov} and in Krattenthaler's
determinant formula \cite[Theorem~1]{kratt}.
Here the denominator gains a new pole with each increase in \(K\).

\section{An integer polynomial in the base}\label{sec:integer}
To compare the evaluations at all conjugates of \(b\), we clear
denominators before specializing \(t\). Let
\(C_K(t)=\prod_{d\le4K}\Phi_d(t)^{E(K,d)}\), where \(\Phi_d\) is the
\(d\)-th cyclotomic polynomial and Appendix~\ref{sec:exponents} lists the
exponents. The arithmetic input from \cite[(4.1), Proposition~4.1
and the proof of Proposition~5.1]{uplift} is
\begin{align}
 C_K\Delta_K&\in\Z[t,t^{-1}][X],\label{eq:laurent}\\
 \deg C_K&=(C+o(1))K^3,\quad
 C=\frac16+\frac{1309}{150\pi^2},\qquad
 \sum_{d\le4K}E(K,d)=O\!\left(K^2\log(2K)\right).             \label{eq:degree}
\end{align}
In units of \(K^3\log b\), Section~\ref{sec:real} gives an upper exponent
of \(-1/3\) for \(\Delta_K(b,\xi)\). Clearing adds \(C\approx1.05086\).
Extracting a common power \(t^{c_K}\) saves \(5/6\), giving the net exponent
\(-1/3+C-5/6\approx-0.11580\). With \(c_K\) given by \eqref{eq:base-exponent},
we use the polynomial
\begin{equation}\label{eq:normalized}
 P_K(t,X)=t^{-c_K}C_K(t)\Delta_K(t,X)\in\Z[t,X],
 \qquad c_K=\tfrac56K^3+O(K^2).
\end{equation}

To verify integrality at \(t=0\), observe that \eqref{eq:laurent} makes
each coefficient in \(X\) an integer Laurent polynomial. The base-power
divisibility in \cite[Proposition~3.1 and
(5.2)]{uplift} gives \(P_K(p,X)\in\Z[X]\) for every positive rational
prime \(p\). If any coefficient had a lowest term \(a t^{-h}\) with
\(h>0\), multiplying its value at any prime \(p\nmid a\) by \(p^h\)
would give an integer congruent to \(a\bmod p\). Integrality of the
value also makes that integer divisible by \(p\), a contradiction.
For \(K=1\) the result reads \(P_1(t,X)=t-(t-1)^2X\).

\section{Decay and nonvanishing at the real base}\label{sec:real}
Specialize \(t=b,\ X=\xi\). The nodes on which \(L\) acts are
\(q,q^3,q^5,\ldots\). A first derivative vanishes at a double zero,
so a polynomial can remove the early, largest nodes from the sum.
For example, \((x-q)^2(x-q^3)^2\) removes the first two. At every
later node its modulus is at most \(q^{2+6}=q^8\).
This is the repeated-root basis used in \cite[\S6]{uplift}. The cubic
gain from a geometric lattice has a fixed-moment antecedent in
Zudilin's estimate \cite[\S4, (8)]{zudilin}.

Take the monic polynomials
\[
 B_j(x)=\prod_{\ell=1}^j\big(x-q^{2\lceil\ell/2\rceil-1}\big),\qquad B_0=1.
\]
Their root exponents sum to \(A_j=\lceil j^2/2\rceil\). Write
\(R_j=2\lceil j/2\rceil-1\) for the largest root exponent for \(j\ge1\),
and set \(R_0=0\). Replacing monomials of degree \(j\) by \(B_j\) in
both rows and columns of \eqref{eq:delta} preserves the determinant,
as both basis changes are triangular with diagonal entries one. Put
\(W=1/D_K\) at \(t=b\). On \([0,q]\), \(W\) and \(W'\) are bounded
uniformly in \(K\). For \(W\) use \(\prod_{a\ge1}(1-q^{2a+1})^{-1}\),
and for \(W'/W\) use \(\sum_{a\ge1}q^{2a}/(1-q^{2a+1})\).

Let \(R=\max(R_i,R_j)\). At each odd \(m<R\) the product \(B_i B_j W\)
has a double zero at \(q^m\). At odd \(m\ge R\) each root factor satisfies
\(\lvert q^m-q^h\rvert\le q^h\). In the Euler derivative \(x\,d/dx\)
a differentiated factor contributes \(q^m\) in place of \(q^h\), gaining
at least the factor \(q^{m-R}\) in the upper bound. A derivative on
\(W\) contributes \(O_q(q^m)\). Summing the resulting geometric bounds
also covers \(i=j=0\), and gives
\[
 \left|L(B_iB_jW)\right|\ll_q (i+j+1)q^{A_i+A_j}.
\]
Extract \(q^{A_i}\) from row \(i\) and column \(i\). Since
\(2\sum_{i<K}A_i=\sum_{i<K}i^2+\lfloor K/2\rfloor\), Hadamard's
inequality gives
\[
  \log|\Delta_K(b,\xi)|\le -\tfrac13 K^3\log b+O_b(K^2).
\]
We take \(\log 0=-\infty\). For fixed \(b>1\), Möbius inversion gives
\(\log|\Phi_d(b)|=\phi(d)\log b+O_b(1)\), where \(\phi\) is Euler's
totient. The error is bounded uniformly in \(d\) by
\(\sum_{h\ge1}|\log(1-b^{-h})|\). Equations \eqref{eq:degree} and
\eqref{eq:normalized} therefore give
\begin{equation}\label{eq:small}
 \log|P_K(b,\xi)|\le(-\delta+o(1))K^3\log b,\qquad
 \delta=\tfrac13+\tfrac56-C=1-\frac{1309}{150\pi^2}>0 .
\end{equation}
Even \(\pi>3\) gives \(\delta>41/1350\).

An upper bound allows the determinant to vanish. Since \(L\) has both
signs, we get nonvanishing by interpolation. In the classical
fixed-moment problem, eventual vanishing of Hankel determinants implies
rationality of their generating function \cite[Arts.~VII and IX]{kron},
\cite[\S4]{zudilin}. Here the denominator \(D_K\) changes with the size.
We give the moving-pole argument of \cite[Proposition~7.1]{uplift} at
the real base \(b\).

\newpage
\begin{lemma}\label{lem:nonzero}
There are arbitrarily large \(K\) for which \(P_K(b,\xi)\ne0\).
\end{lemma}
\begin{proof}
The signed Cauchy transform
\[
 \FF(z)=L((z-x)^{-1})=
 \sum_{\substack{m\ge1\\m\ {\rm odd}}}\frac{\chi(m)q^m}{(z-q^m)^2}
\]
is analytic at infinity and meromorphic on \(\C\setminus\{0\}\).
The series converges normally away from the nodes and their limit point.
Every interior node \(q,q^3,\ldots\) contributes a nonzero double pole,
so \(\FF\) cannot be rational. If the determinants vanish eventually,
the interpolation errors at successive sizes will share zeros at the
exterior nodes \(b^2,b^4,\ldots\), before any common factors are cancelled.

Suppose the determinants are singular at all sufficiently large sizes
\(k\). A kernel vector yields a nonzero polynomial \(Q_k\), of degree
less than \(k\), with \(L(x^j Q_k/D_k)=0\) for \(0\le j<k\).
We use these \(k\) equations to build an approximation to \(\FF\) whose
error vanishes at the exterior nodes shared with the next size. Set
\[
 E_k(z)=D_k(z)L\!\left(\frac{Q_k(x)}{(z-x)D_k(x)}\right),\qquad
 T_k(z)=\FF(z)Q_k(z)-E_k(z).
\]
The polynomial \(Q_k(z)D_k(x)-D_k(z)Q_k(x)\) is divisible by \(z-x\),
so \(T_k\) is a polynomial of degree at most \(k-1\). In the expansion
of \(1/(z-x)\) at infinity the first \(k\) terms give zero after
applying \(L\) to their products with \(Q_k/D_k\). The geometric series
converges in \(C^1[0,q]\) for \(|z|>q\). Multiplication by the degree-\(k\)
polynomial \(D_k(z)\) thus leaves \(E_k(z)=O_k(1/z)\).
Also \(E_k(b^{2a})=0\) for \(1\le a\le k\), because the inner transform
is regular there.

Adjacent sizes share the first \(k\) exterior nodes. Consequently the
polynomial
\[
 T_{k+1}Q_k-T_k Q_{k+1}=E_k Q_{k+1}-E_{k+1}Q_k
\]
vanishes at these \(k\) nodes and has degree at most \(k-1\) by its
growth at infinity. It is zero. Thus all sufficiently large \(k\)
give one reduced fraction \(T_k/Q_k=T/Q\). Write \(Q_k=V_kQ\) and
\(T_k=V_kT\), taking \(Q=1\) if \(T=0\). An exterior interpolation
condition can disappear on reduction where the common factor \(V_k\)
vanishes. The order of the error at infinity bounds how many such points
there can be.

The function \(E=\FF Q-T\) is nonzero since \(\FF\) is nonrational.
It is meromorphic at infinity and \(V_k E=E_k=O_k(1/z)\). Hence \(E\)
has a zero at infinity of some finite order \(h\ge1\), and
\(\deg V_k\le h-1\). In particular \(E(b^{2a})\ne0\) for all
sufficiently large \(a\). Each such \(a\le k\) gives a distinct zero
of \(V_k\), eventually exceeding its degree. This contradiction proves
nonvanishing of \(\Delta_K(b,\xi)\) along an unbounded sequence.
The factors \(b^{-c_K}\) and \(C_K(b)\) are nonzero.
\end{proof}

\newpage
\section{The conjugate estimate and the norm}\label{sec:norm}
At a conjugate base \(t\) with \(|t|\le1\), the second argument of
\(P_K\) could be any fixed complex number. Bounding the factors of
\eqref{eq:normalized} separately can lose the cubic decay. For instance,
at the Lucas conjugate \(t=\varphi^{-2}\), the extracted factor
\(t^{-c_K}=\varphi^{2c_K}\) alone grows on the cubic scale.
Near the unit circle there is also the problem that the denominators in
\eqref{eq:moments} approach zero at roots of unity. Following
\cite[\S4]{weight3}, we estimate on the circle of radius \(1+1/K\).
Here the denominators are separated from zero, and raising the radius
to a power of order \(K^3\) costs only \(\exp(O(K^2))\).
The maximum modulus principle then bounds the whole polynomial on the
closed unit disk, including conjugates of modulus one.

\Needspace{7\baselineskip}
\begin{lemma}\label{lem:disk}
For every fixed \(Y\in\C\),
\[
 \max_{|t|\le1}\logp|P_K(t,Y)|=O_Y(K^{5/2}\log^2(2K))=o(K^3),
 \qquad \logp v=\log\max(1,v).
\]
\end{lemma}
\begin{proof}
Write \(s=1+1/K\) and let \(|t|=s\). Then \(1<|r_a|\le e^2\).
For \(a\ne c\), \(|1-r_a/r_c|\ge1-s^{-1}=1/(K+1)\).
The partial-fraction coefficients of \(1/D_K\) are
\(w_a=\prod_{c\ne a}(1-r_a/r_c)^{-1}\), and
\(|w_a|\le(K+1)^{K-1}\). For positive \(m,j\) the reverse triangle
inequality gives \(|t^m-1|,\ |1+t^{2j}|\ge1/K\).
For \(1\le a\le K,\ 1\le j\le2K\), formula \eqref{eq:moments}
therefore gives \(S_{2a}=O(K^3)\) and \(\mu_j=O(K^2)\).

To bound entries including their polynomial quotients, let
\(U_{a,k}=\LL_{t,Y}(x^k/(1-x/r_a))\). Polynomial division gives
\[
  U_{a,0}=(-1)^a(Y-S_{2a}),\qquad U_{a,k+1}=r_a(U_{a,k}-\mu_k).
\]
Iteration for \(0\le k\le2K-2\) bounds \(|U_{a,k}|\) by
\(e^{4K}(|Y|+O(K^3))\). Each entry in \eqref{eq:delta} is
\(\sum_a w_aU_{a,i+j}\), so determinant expansion gives
\(\logp|\Delta_K(t,Y)|=O_Y(K^2\log(2K))\).

The factors \(1-t^{-h}\) in the Möbius product for \(\Phi_d(t)\) have
modulus between \(1/(K+1)\) and \(2\). Thus
\[
 \log|\Phi_d(t)|\le\phi(d)\log s+\tau(d)\log(2K+2),
\]
where the divisor count satisfies \(\tau(d)\le2\sqrt d\le4\sqrt K\)
for \(d\le4K\). Summing with weights \(E(K,d)\) and using
\eqref{eq:degree} gives \(\logp|C_K(t)|=O(K^{5/2}\log^2(2K))\).
Finally \(|t^{-c_K}|\le1\) on this circle. The maximum modulus principle
applied to the polynomial \(P_K(\,\cdot\,,Y)\) proves the assertion.
\end{proof}

\begin{proof}[Proof of Theorem~\ref{thm:bases}]
Assume that \(\xi=\beta_{1/b}(2)\) belongs to \(F=\Q(b)\), and take a
positive integer \(d\) such that \(d\xi\) is an algebraic integer.
Since \(b\) is an algebraic integer, \(P_K\in\Z[t,X]\), and
\(\deg_X P_K\le K\), the number \(y_K=d^KP_K(b,\xi)\) is an algebraic
integer. By Lemma~\ref{lem:nonzero}, \(y_K\ne0\) for arbitrarily large
\(K\). At those sizes the absolute norm is a positive integer.

Let \(\iota\) be the chosen real embedding of \(F\). The other embeddings
\(\sigma:F\hookrightarrow\C\) satisfy \(\lvert\sigma(b)\rvert\le1\).
For each of them \(Y=\sigma(\xi)\) is fixed as \(K\) grows, so
Lemma~\ref{lem:disk} bounds its contribution to the logarithmic norm by
\(o(K^3)\). There are only \([F:\Q]-1\) such embeddings. Using
\eqref{eq:small} at \(\iota\) gives
\begin{align*}
 0\le\log|N_{F/\Q}(y_K)|
 &= [F:\Q]\,K\log d+\log|P_K(b,\xi)|
       +\sum_{\sigma\ne\iota}\log|P_K(\sigma(b),\sigma(\xi))|\\
 &\le -\delta(\log b)K^3+o(K^3).
\end{align*}
Since \(b>1\) and \(\delta>0\), the right-hand side tends to minus infinity.
\end{proof}

\clearpage
\appendix
\section{The cyclotomic factors and the power of the base}\label{sec:exponents}
For clarity, here is the normalization behind \eqref{eq:normalized}.
Write \(v_+=\max(v,0)\). Define \(f(0)=0\) and
\(f(m)=\lfloor(m+2)^2/4\rfloor\) for \(m\ge1\).
In the first case below write \(K=rd+s,\ 0\le s<d\).
The exponents of \cite[(4.1)]{uplift} are
\[
 E(K,d)=
 \begin{cases}
 s f(r+1)+(d-s)f(r),&d\text{ odd},\ d\le K,\\
 2(K-\lfloor(d-1)/2\rfloor),&d\text{ odd},\ K<d\le2K,\\
 0,&d\text{ odd},\ d>2K,\\
 \displaystyle\sum_{j\ge1,\ j\ {\rm odd}}(K-jh)_+,&d=2h,\ h\text{ odd},\\
 \min\{K,(K-4u)_++(K-u-1)_+\},&d=4u .
 \end{cases}
\]
The last line illustrates how the signs affect the arithmetic. At a
primitive root of order \(4u\), the signed proper-fraction residue has
rank at most \((K-4u)_+\) because the interpolated numerator has a
double zero at each cluster of coincident nodes. The polynomial quotient
adds rank at most \((K-u-1)_+\). These are the two terms in the pole-order
bound of \cite[Proposition~4.1]{uplift}.

For a direct calculation of the power of \(t\), let \(\nu=\ord_{t=0}\).
In the formal series ring set \(Z_0=\sum_{m\ge1}\chi(m)g(t^m)\) and
\(H(x)=\sum_{j\ge1}\mu_j(t)x^j\).
Expanding \(g\) gives \(H(x)=\sum_{m\ge1}\chi(m)g(t^m x)\). The same
shift as in Section~\ref{sec:pencil} gives
\(S_{2a}(t)=Z_0-(-1)^aH(r_a)\). With \(U=X-Z_0\), the pole value of
\(\LL_{t,X}\) is \((-1)^a U+H(r_a)\).

Put \(F_K(x)=\prod_a(x-r_a)\). In \(\LL_{t,X}(x^k/F_K(x))\) the
coefficient of \(U\) equals
\[
      -\sum_{a=1}^K\frac{(-1)^a r_a^{k-1}}{F'_K(r_a)} .
\]
The constant part belongs to \(\Z[[t]]\). Indeed the quotient by the
monic \(F_K\) has integral coefficients and moments. The proper part
contributes \(-\sum_{j\ge1}\mu_j(t)h_{k+j-K}(r_1,\ldots,r_K)\),
where \(h_\ell\) is the complete homogeneous symmetric polynomial,
\(h_0=1\), and \(h_\ell=0\) for \(\ell<0\).
To see this, expand \(1/F_K\) at infinity and in simple fractions.
Comparing coefficients gives
\(\sum_a r_a^\ell/F'_K(r_a)=h_{\ell-K+1}\) for \(\ell\ge0\).
Here \(\ell=k+j-1\ge0\), and \(\mu_j(t)\in t^j\Z[[t]]\) ensures
convergence.

The equality \(\nu(r_a-r_h)=2\min(a,h)\) for \(a\ne h\) now gives
\[
 \nu\!\left(\frac{r_a^{k-1}}{F'_K(r_a)}\right)
 =a^2+(2k-2K-1)a\ge -w(k),\qquad
 w(k)=(K-k)_+\big((K-k)_++1\big).
\]
The inequality follows by minimizing the quadratic over integers
\(1\le a\le K\). Translating \(U\) back to \(X\) preserves the bound.
For a determinant term indexed by a permutation \(\sigma\), convexity
of \(w\) gives \(2w(i+\sigma(i))\le w(2i)+w(2\sigma(i))\). Thus
\(\sum_i w(i+\sigma(i))\le\sum_i w(2i)\). Restoring the factor
\((-1)^K t^{K(K+1)}\) in \(1/D_K=(-1)^K t^{K(K+1)}/F_K\) in every row
gives the coefficientwise bound \(\nu(\Delta_K)\ge c_K\) with
\begin{equation}\label{eq:base-exponent}
 c_K=K^2(K+1)-\sum_{i=0}^{K-1} w(2i)
    =\frac{10K^3+3K^2-10K-3(K\bmod 2)}{12}.
\end{equation}
The last equality follows by summing separately for even and odd \(K\).

\section*{Verification}
Two collaborating AI models reviewed the mathematics and exposition of the
initial 18 November 2026 draft and reproduced its three-pole calculation,
taking the cited cyclotomic estimates as inputs. A separate AI model
audited that draft for exposition, and a collaborating model checked the
algebraic-base extension in this revision.

\clearpage
\small
\begin{thebibliography}{10}
\setlength{\itemsep}{2pt}
\bibitem{dnns}
D. Duverney, Ke. Nishioka, Ku. Nishioka and I. Shiokawa,
Transcendence of Rogers--Ramanujan continued fraction and reciprocal
sums of Fibonacci numbers, \emph{Proc. Japan Acad. Ser. A Math. Sci.}
\textbf{73} (1997), 140--142.
\href{https://doi.org/10.3792/pjaa.73.140}{doi:10.3792/pjaa.73.140}.
See also \emph{RIMS K\^oky\^uroku} \textbf{1060} (1998), 91--100.
\href{https://www.kurims.kyoto-u.ac.jp/~kyodo/kokyuroku/contents/pdf/1060-10.pdf}{Kyoto account}.
\bibitem{jm}
F. Jouhet and E. Mosaki,
Diophantine properties for \(q\)-analogues of Dirichlet's beta function
at positive integers, \emph{Trans. Amer. Math. Soc.} \textbf{363}
(2011), 1533--1554. \href{https://doi.org/10.1090/S0002-9947-2010-05236-5}{doi:10.1090/S0002-9947-2010-05236-5}.
\bibitem{bz}
P. Bundschuh and W. Zudilin, Irrationality measures for certain
\(q\)-mathematical constants, \emph{Math. Scand.} \textbf{101}
(2007), 104--122.
\href{https://doi.org/10.7146/math.scand.a-15034}{doi:10.7146/math.scand.a-15034}.
\bibitem{uplift}
uplift-universal, \emph{Irrationality of the \(q\)-Catalan value
\(\beta_{1/p}(2)\)}, October 2026.
\href{https://github.com/uplift-universal-inc/q-catalan-irrationality/blob/main/paper/main.pdf}{Manuscript}.
\bibitem{wang}
W. Wang, Evaluation of reciprocal sums of hyperbolic functions using
quasimodular forms, \emph{J. Math. Anal. Appl.} \textbf{549} (2025),
no.~2, 129526. \href{https://doi.org/10.1016/j.jmaa.2025.129526}{doi:10.1016/j.jmaa.2025.129526}.
\href{https://arxiv.org/abs/2312.16929v1}{Preprint version}.
\bibitem{weight3}
twinheterodox, \emph{The weight-three Lambert sum at the golden ratio is
irrational}, 2026. \href{https://hexagonmath.org/2611.00623v1}{Hexagon 2611.00623v1}.
\bibitem{krisman}
S. Krisman, \emph{The sum of the reciprocal cubes of the Fibonacci
numbers is irrational}, 2026. \href{https://hexagonmath.org/2611.00473v1}{Hexagon 2611.00473v1}.
\bibitem{bv}
P. Bundschuh and K. Väänänen, Linear independence of certain Lambert
and allied series, \emph{Acta Arith.} \textbf{120} (2005), 197--209.
\href{https://doi.org/10.4064/aa120-2-7}{doi:10.4064/aa120-2-7}.
\bibitem{uvarov}
V. B. Uvarov, The connection between systems of polynomials orthogonal
with respect to different distribution functions,
\emph{USSR Comput. Math. Math. Phys.} \textbf{9} (1969), no.~6, 25--36.
\href{https://doi.org/10.1016/0041-5553(69)90124-4}{doi:10.1016/0041-5553(69)90124-4}.
\bibitem{kratt}
C. Krattenthaler, A determinant identity for moments of orthogonal
polynomials that implies Uvarov's formula for the orthogonal polynomials
of rationally related densities, 2021.
\href{https://arxiv.org/abs/2103.03969}{arXiv:2103.03969}.
\bibitem{zudilin}
W. Zudilin, On the irrationality of generalized \(q\)-logarithm,
\emph{Res. Number Theory} \textbf{2} (2016), article 15.
\href{https://doi.org/10.1007/s40993-016-0042-x}{doi:10.1007/s40993-016-0042-x}.
\bibitem{kron}
L. Kronecker, Zur Theorie der Elimination einer Variabeln aus zwei
algebraischen Gleichungen, \emph{Monatsber. K\"onigl. Preuss. Akad.
Wiss. Berlin} (1881), 535--600.
\href{https://www.e-rara.ch/download/pdf/5761162.pdf\#page=42}{Collected works, vol.~II}.
\end{thebibliography}
\end{document}

